\subsection{Geometric Interpretation}
\label{sec:geometric-interpretation}
\slidesrc{Distributional semantics: 22--27; Semantics (long version): 6}

\subsubsection{Words as points}
\slidesrc{Distributional semantics: 22--24}

The row vector $\mathbf{x}_{\text{dog}} = (115, 83, 10, 42, 33, 17)$ describes how
\emph{dog} is used in the corpus. It can be seen as the coordinates of a \textbf{point}
in an $n$-dimensional Euclidean space ($n = 6$ here, one dimension per context). We can
draw two dimensions, \emph{get} and \emph{use}: $\mathbf{x}_{\text{dog}} = (115, 10)$,
$\mathbf{x}_{\text{cat}} = (52, 4)$, $\mathbf{x}_{\text{knife}} = (51, 84)$,
$\mathbf{x}_{\text{boat}} = (59, 23)$.

\begin{center}
\begin{tikzpicture}
\begin{axis}[width=10cm, axis equal image, xmin=0, xmax=130, ymin=0, ymax=105,
    xlabel={get}, ylabel={use}, title={Two dimensions of the English V--Obj DSM},
    title style={font=\small}, tick label style={font=\scriptsize}, label style={font=\small},
    axis lines=left, clip=false]
  % unit-length arc (radius 80, as in the slides)
  \addplot[domain=0:90, samples=60, nbgray, thin] ({80*cos(x)}, {80*sin(x)});
  \draw[-{Stealth}, thick, nbgreen]  (axis cs:0,0) -- (axis cs:51,84);
  \draw[-{Stealth}, thick, nbblue]   (axis cs:0,0) -- (axis cs:59,23);
  \draw[-{Stealth}, thick, black]    (axis cs:0,0) -- (axis cs:52,4);
  \draw[-{Stealth}, thick, nbred]    (axis cs:0,0) -- (axis cs:115,10);
  \fill[nbgreen] (axis cs:51,84) circle (2pt) node[above left, font=\small] {knife};
  \fill[nbblue]  (axis cs:59,23) circle (2pt) node[above, font=\small] {boat};
  \fill[black]   (axis cs:52,4)  circle (2pt) node[above left, font=\small] {cat};
  \fill[nbred]   (axis cs:115,10) circle (2pt) node[above, font=\small] {dog};
  % normalised points on the arc
  \draw[nbgreen, thick] (axis cs:41.5,68.4) circle (2.5pt);
  \draw[nbblue, thick]  (axis cs:74.5,29.1) circle (2.5pt);
  \draw[nbred, thick]   (axis cs:79.7,7.0)  circle (2.5pt);
  % angle between knife and dog
  \addplot[domain=5:58.7, samples=40, dashed, thick] ({48*cos(x)}, {48*sin(x)});
  \node[font=\small] at (axis cs:48,50) {$\alpha$};
\end{axis}
\end{tikzpicture}
\captionof{figure}{Words as vectors in the (get, use) plane. Circles on the arc: the same
vectors normalised to the same length (only the direction is kept). \emph{Dog} and
\emph{cat} point in almost the same direction even though they are far apart as points.}
\end{center}

A first idea: \textbf{similarity = spatial proximity}, measured with the Euclidean
distance. But the location of a point depends on the \textbf{frequency} of the noun:
\emph{dog} is about 2.7 times more frequent than \emph{cat}, so all its counts are larger
and its point lies much farther from the origin.

\begin{nbexample}{Euclidean distance is misleading}
\begin{align*}
  d(\text{dog}, \text{cat})  &= \sqrt{(115-52)^2 + (10-4)^2}  = \sqrt{3969 + 36}  \approx 63.3\\
  d(\text{dog}, \text{boat}) &= \sqrt{(115-59)^2 + (10-23)^2} = \sqrt{3136 + 169} \approx 57.5
\end{align*}
By Euclidean distance, \emph{boat} is closer to \emph{dog} than \emph{cat} is: a
consequence of frequency, not of meaning.
\end{nbexample}

\subsubsection{Direction, not location}
\slidesrc{Distributional semantics: 25--27}

A vector can also be seen as an \textbf{arrow from the origin}. Its \emph{length} mostly
reflects how frequent the word is; its \emph{direction} reflects \emph{the proportions}
in which the word occurs with the different contexts, i.e.\ how it is used. So the
direction is more important than the location. Two equivalent solutions:
\begin{itemize}[nosep]
  \item use the \textbf{angle} $\alpha$ between two vectors as a distance measure;
  \item or \textbf{normalise} the length $\norm{\mathbf{x}}$ of every arrow to the same
        value, and then compare points (the circles on the arc in the figure).
\end{itemize}

\begin{nbdef}{Cosine similarity}
The \textbf{cosine similarity} of two vectors is the cosine of the angle between them:
\[
  \cos\alpha = \frac{\mathbf{x}\cdot\mathbf{y}}{\norm{\mathbf{x}}\,\norm{\mathbf{y}}}
             = \frac{\sum_i x_i y_i}{\sqrt{\sum_i x_i^2}\,\sqrt{\sum_i y_i^2}}.
\]
It is 1 for vectors pointing in the same direction, 0 for orthogonal vectors (with
non-negative counts it is never negative), and it does not change if a vector is
multiplied by a positive number: it ignores frequency.
\end{nbdef}

\begin{nbexample}{dog vs.\ cat and dog vs.\ knife in two dimensions}
\begin{align*}
  \cos(\text{dog}, \text{cat}) &= \frac{115\cdot 52 + 10\cdot 4}{\sqrt{115^2+10^2}\,\sqrt{52^2+4^2}}
     = \frac{6020}{115.43 \cdot 52.15} \approx 1.000 \quad(\alpha \approx 0.6^\circ)\\
  \cos(\text{dog}, \text{knife}) &= \frac{115\cdot 51 + 10\cdot 84}{115.43 \cdot 98.27}
     = \frac{6705}{11343} \approx 0.591 \quad(\alpha \approx 53.8^\circ)
\end{align*}
Now \emph{cat} is (in these two dimensions) practically identical to \emph{dog}, and
\emph{knife} is far. The slides report $\alpha = 54.3^\circ$ for knife--dog; computing it
from the counts gives $53.8^\circ$, the small difference does not matter.
\end{nbexample}

\begin{nbwhy}
Every modern word representation, from count-based DSMs to \textbf{word embeddings}
(\emph{word2vec}, GloVe) to the internal vectors of BERT and GPT, follows this recipe:
a word is a vector learnt from its contexts, and similar words have vectors with a small
angle between them. Word2vec learns short dense vectors (a few hundred dimensions)
directly, instead of counting; implementations in Python and C, with examples, are linked
in the course's AulaWeb module.
\end{nbwhy}

\begin{nbsummary}[distributional semantics]
\begin{itemize}[nosep]
  \item Wittgenstein: meaning = use; Firth: ``you shall know a word by the company it
        keeps''; Harris: difference of meaning correlates with difference of
        distribution.
  \item \emph{bardiwac}: the meaning of an unknown word can be inferred from its
        contexts alone.
  \item Co-occurrence matrix: rows = target words, columns = contexts, cells = counts;
        each row is the word's vector. Weight the counts (e.g.\ $\log(1+f)$) before
        comparing.
  \item Euclidean distance is dominated by frequency; use the \textbf{angle}, i.e.\
        \textbf{cosine similarity} $\mathbf{x}\cdot\mathbf{y}/(\norm{\mathbf{x}}\norm{\mathbf{y}})$,
        or normalise the vectors.
  \item This is the basis of word embeddings (word2vec) and of the representations in
        LLMs.
\end{itemize}
\end{nbsummary}
